% Generated draft: requires mathematical and source review.
\documentclass[11pt]{article}
\usepackage[margin=25mm]{geometry}
\usepackage{fontspec}
\setmainfont{Latin Modern Roman}
\usepackage{amsmath,amssymb,mathtools,unicode-math}
\setmathfont{Latin Modern Math}
\usepackage{enumitem,needspace,xurl,hyperref}
\setlength{\parindent}{0pt}
\setlength{\parskip}{5pt}
\setlength{\emergencystretch}{3em}
\newcommand{\N}{\mathbb{N}}
\newcommand{\Z}{\mathbb{Z}}
\newcommand{\Q}{\mathbb{Q}}
\newcommand{\R}{\mathbb{R}}
\newcommand{\C}{\mathbb{C}}
\newcommand{\F}{\mathbb{F}}
\newcommand{\A}{\mathbb{A}}
\newcommand{\PP}{\mathbb{P}}
\newcommand{\E}{\mathbb{E}}
\newcommand{\eps}{\varepsilon}
\newcommand{\dd}{\,\mathrm{d}}
\begin{document}
\section*{UnsolvedMath ID: 30002780}
\textbf{Source problem code:} catalogue entry
\subsection{Definitions and mathematical statement}
Let \(K\) denote the triangulation and let \(d=\dim K\), as in the supplied record. The record describes \(K\) as a triangulated homology manifold, but does not specify homology coefficients, connectedness, orientability, compactness, or whether boundary is allowed.

For reference, if \(K\) is a finite \(d\)-dimensional simplicial complex, write \(f_r(K)\) for the number of its \(r\)-dimensional faces and set \(f_{-1}(K)=1\). The usual face-number convention defines \(h_0(K),\ldots,h_{d+1}(K)\) by
\[
\sum_{r=0}^{d+1}h_r(K)t^{d+1-r}
=\sum_{r=0}^{d+1}f_{r-1}(K)(t-1)^{d+1-r},
\]
and defines \(g_i(K)=h_i(K)-h_{i-1}(K)\) for \(1\leq i\leq d+1\). These are standard candidate conventions for interpreting the supplied notation; the record does not explicitly give its definitions of \(g_i\) or confirm finiteness. It describes \(\mu_j(K)\) only as the averaged discrete Morse numbers of the triangulation. The averaging procedure, normalization, and coefficient dependence are not supplied, so this description does not uniquely define \(\mu_j(K)\).

\textbf{Recorded conjectural inequality (attributed to Bagchi).} For every intended admissible \(d\)-dimensional triangulated homology manifold \(K\), and every intended relevant index \(i\), the proposed bound is
\[
 g_i(K)\geq \binom{d+1}{i}
 \left((-1)^i+\sum_{j=1}^{i}(-1)^{i-j}\mu_{j-1}(K)\right).
\]
The displayed dimension parameter and binomial coefficient are retained exactly as supplied. The admissible manifold class, the index range, and the precise \(\mu\)-vector definition remain unresolved; consequently, this is a statement template requiring source clarification, rather than a fully specified universally quantified conjecture.
\subsection{Short English statement}
The record attributes to Bagchi a lower bound on the \(g\)-numbers of triangulated homology manifolds in terms of alternating sums of their \(\mu\)-numbers. The intended manifold hypotheses, index range, and definition of the \(\mu\)-numbers must be clarified before the conjecture has a precise formulation.
\subsection{Sources}
\begin{itemize}
\item Local UnsolvedMath registry: \texttt{data/unsolvedmath/problems.json}, permanent ID \(30002780\), title ``Bagchi's Lower Bound for Triangulated Manifolds,'' source problem code \texttt{catalogue entry}. The record attributes the conjecture to Bagchi and gives proposed year 2015; no separate proposer is specified.
\item Supplied report citation: \emph{Geometric and Algebraic Combinatorics} (2015), Oberwolfach Reports, pp.\ 285--368, DOI \texttt{10.4171/owr/2015/5}. \url{https://doi.org/10.4171/owr/2015/5}. The supplied background describes these reports as documenting mathematical workshops and their open-problem discussions.
\item Supplied literature reference, labeled ``Bagchi, the mu vector and generalized lower bounds'': \url{https://doi.org/10.1016/j.ejc.2015.04.007}.
\item Supplied literature reference, labeled ``original OWR report'': \url{https://doi.org/10.4171/owr/2015/5}.
\item Literature assessment dated exactly \texttt{2026-08-21}: the supplied curation labels the problem partially solved and states that the \(\mu\)-vector lower-bound conjecture is proved for connected closed combinatorial \(3\)-manifolds and for broader locally tame classes, while the general homology-manifold statement in the record remains unproved. This is supplied dated curation, not an independently verified assessment. The background requests independent verification of time-sensitive status and cited claims.
\item The supplied statement-status label is \texttt{corrected\_verified}. Its verification note reports that the official Oberwolfach passage was checked directly and that the question, definitions, and notation were restored without strengthening the claim. This draft preserves that provenance assertion without claiming an independent source visit or verification.
\end{itemize}
\end{document}
